% =====================================================================
%  APPENDIX: Self-referential ignorance and Hardy's continuity axiom
%  -- the owed arrow, its decomposition, and open problems
%  ADDED 2026-07-17. Consolidates the session records
%  continuity_limit_realisability_2026-07-04.md and
%  continuity_blocks_and_routes_2026-07-17.md into paper form.
%  Requires (amsthm): \newtheorem{lemma}{Lemma}, {theorem}, {corollary},
%  {remark}, and \newtheorem{openproblem}{Open Problem} -- the last
%  matching the numbering stream of op:dictionary-lemmas if shared.
%  All \cite keys are stubs; a commented bibliography list is at the
%  end of this file. VERIFY flags mark memory-based citations.
% =====================================================================
%
% ---------------------------------------------------------------------
%  MAIN-BODY INSERT (place in sec:self-reference at the continuity
%  passage; adjust the lead-in to the surrounding prose):
%
%  It would complete the arch of this section to prove that
%  self-referential ignorance \emph{implies} continuity in Hardy's
%  form -- his Axiom~5, a continuous reversible transformation between
%  any two pure states -- so that the reconstruction theorems could be
%  cited with no axiom left unexplained. A conclusive proof remains
%  elusive. The obstruction, however, is no longer vague: the
%  derivation decomposes into a chain in which every link but three is
%  a theorem, and the three that remain are stated as sharply
%  constrained open problems in Appendix~\ref{app:continuity-open}.
%  That appendix also records why the tempting short routes fail
%  (uncountability is not continuity; completeness is not
%  connectedness; superposition alone does not force a continuum),
%  which established mathematics bears on each link (automatic
%  continuity, the classification of the completions of $\mathbb{Q}$,
%  Specker sequences), and what would count as a proof.
% ---------------------------------------------------------------------

\section{Self-referential ignorance and Hardy's continuity axiom: status and open problems}
\label{app:continuity-open}

This appendix isolates the one arrow of the reconstruction this paper does not prove:
that self-referential ignorance implies Hardy's continuity axiom \cite{hardy2001}
(\emph{Axiom~5}: ``there exists a continuous reversible transformation on a system between
any two pure states of that system''). The strategy of \S\ref{sec:self-reference} is to
derive Hardy's premises and then cite his theorem for everything downstream --- complex
Hilbert space, the quadratic norm, $U(n)$. Four of the five premises are treated in the
main text; continuity is the residue. We state here exactly what is proved, exactly what
is assumed, why the naive routes fail, and which sharply bounded open problems separate
the current position from a theorem. The appendix is written to be self-contained, since
we believe the question --- posed at this level of detail --- is new, and we would like it
to be answerable by readers who do not share the framework's other commitments.

\subsection{What ``continuity'' asserts: six inequivalent notions}
\label{app:cont-ladder}

Hardy's axiom does not assert the continuity of a function. It asserts the
\emph{path-connectedness of a group action}: the reversible transformations form a group
$G$ acting on the pure states, and any two pure states lie on a continuous $G$-path. At
least six inequivalent notions circulate in informal discussions of ``continuity,'' and
the failed short routes of \S\ref{app:cont-deadroutes} each conflate two of them:
(i)~topological continuity of a map; (ii)~connectedness of a space, and the strictly
stronger path-connectedness; (iii)~Cauchy completeness of a uniform space (every Cauchy
sequence has a limit); (iv)~order-completeness plus the Archimedean property --- Dedekind's
\emph{Stetigkeit}, which for ordered fields is categorical for $\mathbb{R}$;
(v)~local compactness, the ambient property on which the classical structure theorems
(Cartan, Baire, Haar) run; (vi)~analytic rigidity over $\mathbb{C}$, where the identity
theorem makes a germ determine a function on a \emph{connected} domain.

\begin{remark}[Three separating witnesses]\label{rem:witnesses}
(a)~\emph{Complete $\neq$ connected}: $\mathbb{Q}_p$ is Cauchy-complete and totally
disconnected; its additive group is moreover divisible, so completeness and divisibility
together, without more, do not yield connectedness.
(b)~\emph{Uncountable $\neq$ connected}: the Cantor set.
(c)~\emph{Connected $\neq$ path-connected}: the $p$-adic solenoid, a compact connected
abelian group with a dense proper path-component; a transformation group of solenoid type
would satisfy every naive ``continuity'' demand and still violate Axiom~5 for pairs of
states in different path-components. What excludes (a) at the group level is the
Archimedean side of the fork of \S\ref{app:cont-scalar} (no-small-subgroups, in the sense
of Hilbert's fifth problem \cite{gleason1952, montgomeryzippin}); what excludes (c) is
the Lie structure supplied by Cartan's theorem once the group acts on a finite-dimensional
real state body. Both exclusions therefore \emph{consume} the real scalars, whose
derivation is Layer~1 below.
\end{remark}

\subsection{The two layers, and where hidden imports sit}
\label{app:cont-layers}

The owed arrow is two arrows. \emph{Layer~1 (scalars)}: why the probability scalars are
$\mathbb{R}$ --- Dedekind-complete, Archimedean, connected --- rather than $\mathbb{Q}$,
the computable reals, a hyperreal field, or a $p$-adic completion. \emph{Layer~2 (group)}:
why, over those scalars, the reversible-transformation group is closed, divisible, and
transitive, whence connected and path-connected (Theorem~\ref{thm:conditional} below).
Keeping the layers separate exposes two hidden imports in the standard treatments. First,
Hardy's Axiom~1 already assumes relative frequencies \emph{tend to definite limits} --- a
limit-realisability premise of exactly the kind at issue. Second, Kolmogorov's axiom of
countable additivity is literally his \emph{continuity} axiom, the one de Finetti rejected
as non-operational \cite{kolmogorov1933, definetti}; the present framework takes the
opposite side of that dispute, and holds that self-reference is what forces the
completion de Finetti resisted. Both imports discharge to the same principle (P1 below).

\subsection{Layer 1: which completion? The Accumulation Lemma}
\label{app:cont-scalar}

Probability values enter operationally as frequency ratios: rationals, refined without
bound, never completed by any internal agent. If the scalar set is a completion of
$\mathbb{Q}$, Ostrowski's theorem \cite{ostrowski1916} forces a fork: every nontrivial
absolute value on $\mathbb{Q}$ is equivalent to the real one or to a $p$-adic one. The
$p$-adic tine is not vacant --- Khrennikov's programme lives on it
\cite{khrennikov-padic} --- so the Archimedean choice must be argued, not assumed. The
super-task thread of \S\ref{sec:self-reference} supplies the argument.

\begin{lemma}[Accumulation selects the Archimedean completion]\label{lem:accumulation}
Consider the Thomson schedule $t_n = 1 - 2^{-n}$ \cite{thomson1954}. The schedule
accumulates ($|t_n - 1| \to 0$) if and only if $|2^{-n}| \to 0$. This holds for the real
absolute value; it fails for every other completion of $\mathbb{Q}$ and for every
non-Archimedean ordered field: for $p$ odd, $|2^{-n}|_p = 1$ for all $n$; for $p = 2$,
$|2^{-n}|_2 = 2^{n} \to \infty$; and in an ordered field with an infinitesimal
$\varepsilon$ one has $0 < \varepsilon < 2^{-n}$ for all $n$, so $2^{-n} \not\to 0$ in the
order topology. Consequently a scalar system in which the super-task can so much as
accumulate carries the Archimedean absolute value, and its Cauchy completion is
$\mathbb{R}$ --- uniquely, since $\mathbb{R}$ is the only complete Archimedean ordered
field.
\end{lemma}

The lemma divides the labour exactly. The \emph{Archimedean} property (shared by
$\mathbb{Q}$) is what lets super-tasks accumulate --- the textbook lamp's limit point
$t=1$ is rational. \emph{Completeness} ($\mathbb{R}$ only) is what guarantees a terminal
point for \emph{generic} super-tasks; \S\ref{app:cont-specker} exhibits the witness. Two
further remarks. The uniformity is not conventional: closeness of frequencies means
proportionally few discordant trials --- the order uniformity; $p$-adic closeness of
counts (congruence modulo $p^n$) has no operational reading. And the framework's
phase mechanism is Archimedean-specific: $\mathbb{Q}_p$ and its unit groups are totally
disconnected, so no $p$-adic scalar system carries a continuous interpolating phase, a
$U(1)$, or a winding integer (cf.\ Appendix~\ref{app:winding-number}).

\subsection{Layer 2: the proved machinery}
\label{app:cont-machinery}

Let $K$ be the state body --- granting Layer~1, a compact convex set with nonempty
interior in $\mathbb{R}^{K}$ --- and let $\mathrm{Sym}(K)$ be its (compact) group of
affine automorphisms. Let $G_0 \subseteq \mathrm{Sym}(K)$ be the group of reversible
transformations \emph{available} in the pre-FANOUT regime (composition and inverses of
available transformations are available; reversibility is the regime's second face,
\S\ref{sec:fanout-boundary}).

\begin{lemma}[Closure preserves divisibility]\label{lem:closure-divisible}
If $H$ is a divisible subgroup of a compact metrizable group, then $\overline{H}$ is
divisible.
\end{lemma}

\begin{proof}
Let $x = \lim_k h_k$ with $h_k \in H$, and fix $n$. Write $h_k = r_k^{\,n}$ with
$r_k \in H$. By compactness pass to $r_{k_j} \to r \in \overline{H}$; by continuity of
the power map, $r^{\,n} = \lim_j r_{k_j}^{\,n} = x$. (Dependent choice suffices
throughout; no stronger choice principle is used in this appendix.)
\end{proof}

\begin{lemma}[Divisible and closed implies connected Lie]\label{lem:closedness}
Let $G \le \mathrm{Sym}(K)$ be topologically closed and divisible. Then $G$ is a compact
connected Lie group; in particular it is path-connected.
\end{lemma}

\begin{proof}
Closed in compact is compact; a closed subgroup of the linear group $\mathrm{Sym}(K)$ is
Lie with finitely many components (Cartan). The component group $\pi_0(G) = G/G_0$ is a
finite quotient of a divisible group, hence divisible; a finite divisible group is
trivial (with $|\pi_0| = m$, every $h$ satisfies $h^m = e$, so only $e$ admits an $m$-th
root). Hence $G = G_0$ is connected, and a connected Lie group is path-connected.
\end{proof}

\begin{theorem}[Conditional derivation of Axiom 5]\label{thm:conditional}
Assume Layer~1, and:
\begin{itemize}
\item[\textup{(P1)}] \emph{Descriptive closure.} The theory's transformation group is
$G = \overline{G_0}$, the closure of the available group in $\mathrm{Sym}(K)$.
\item[\textup{(P2$'$)}] \emph{Divisibility.} Every available transformation is realisable
in $n$ equal available stages, for every $n$: $G_0$ is divisible.
\item[\textup{(P3)}] \emph{No superselection.} $G$ acts transitively on the pure states
(extreme points of $K$).
\end{itemize}
Then $G$ is a compact connected Lie group acting transitively on pure states, and for any
two pure states $x, y$ there is a continuous path $t \mapsto g(t)$ of reversible
transformations with $g(0) = e$ and $g(1)x = y$: Hardy's Axiom~5.
\end{theorem}

\begin{proof}
$G$ is closed by (P1) and divisible by (P2$'$) with Lemma~\ref{lem:closure-divisible};
Lemma~\ref{lem:closedness} makes it compact, connected, Lie, path-connected. By (P3)
choose $g$ with $gx = y$ and a path from $e$ to $g$.
\end{proof}

\begin{remark}[A weaker premise buys less]\label{rem:fallback}
Replacing (P2$'$) by the weaker \emph{no minimal step} --- the identity is not isolated in
$G_0$, since a gap around ``do nothing'' would be a smallest detectable change, a
self-clock --- still forces $\overline{G_0}$ to be non-discrete, hence of dimension
$\ge 1$ with nontrivial identity component: continuous reversible paths \emph{exist}. But
$\pi_0$ may survive, and Axiom~5 is then available only within path-components. The full
axiom needs full divisibility. We record the gap rather than paper over it.
\end{remark}

\subsection{The premises: grounding and status}
\label{app:cont-premises}

\emph{(P1), descriptive closure,} is the framework-native premise. Its grounding is the
founding obstruction of \S\ref{sec:self-reference}: no physical system can specify its
own state to unbounded precision (Church--Turing; \cite{popper1950}; Lawvere's
fixed-point form \cite{lawvere1969}). States and transformations are therefore
finite-precision description-classes; a coherent family of ever-finer descriptions is a
Cauchy class; and the space of Cauchy classes is complete \emph{by the universal property
of completion} --- not because the world contains actual infinite limits, but because the
theory's objects are descriptions-up-to-operational-indistinguishability. Three
independent mathematical literatures corroborate the principle by deriving continuity
from ``being given by data'': measurable homomorphisms of Polish groups are continuous
(Steinhaus--Pettis--Christensen; survey \cite{rosendal2009}); every abstract homomorphism
from a compact \emph{semisimple} Lie group into a compact group is continuous, so the
topology of $SU(n)$ is remembered by its algebra (van der Waerden
\cite{vanderwaerden1933}; cf.\ \cite{tsankov2013}); and in computable analysis and domain
theory, \emph{computable implies continuous} is a theorem \cite{weihrauch2000,
abramskyjung, smyth}. The tori are the pointed exception to the second face: $U(1)$ has
discontinuous abstract automorphisms under choice, so the one component of quantum
symmetry whose topology algebra does \emph{not} remember is the phase --- precisely the
component this framework identifies with self-referential ignorance. The piece of
topology that self-reference must personally supply is the piece that encodes it.
(The same Andrew Gleason stands at both ends of this bridge: the no-small-subgroups case
of Hilbert's fifth problem \cite{gleason1952}, and the frame-function theorem, whose
analytic core is again continuity-from-additivity \cite{gleason1957}.)

\emph{(P2$'$), divisibility,} is grounded in the anti-event. An un-happening with no
distinguished intermediate stage must admit an $n$-th stage for every $n$; over
$\mathbb{R}$ the anti-event $-1$ has no stage at all even for $n = 2$ --- $\sqrt{-1}$ is
half an anti-event, and this, not calculational convenience, is where the scalars leave
the real line. The minimal divisible group containing an element of order two is its
injective hull, the Pr\"ufer group $\mathbb{Z}(2^{\infty})$: exactly the dyadic tower
$\{\pm 1, \pm i, e^{\pm i\pi/4}, \dots\}$, i.e.\ the Thomson schedule's halvings read as
phases. Its closure is $U(1)$. Thus the two premises (P2$'$)~then~(P1), applied to the
single primitive $-1$, construct the phase group:
\begin{equation}
U(1) \;=\; \overline{\,\mathrm{divisible\ hull}(-1)\,}
\;=\; \overline{\,\mathbb{Z}(2^{\infty})\,},
\label{eq:dyadic-tower}
\end{equation}
and the same two premises, applied one level up to $G_0$, are the hypotheses of
Lemma~\ref{lem:closedness}. One motor, two layers. What (P2$'$) still lacks is a
\emph{derivation} from the regime's defining property; Open Problem~\ref{op:stage-lattice}
states the candidate reduction.

\emph{(P3), no superselection,} we expect to be a restatement of the regime's third face
(``no classical record''): a superselection sector label would be a basis-free,
measurement-free, absolute distinguished fact --- a classical observable commuting with
everything, i.e.\ a record predating FANOUT. Open Problem~\ref{op:transitivity} asks for
the orbit-level argument.

\subsection{Why each premise is necessary: the foils}
\label{app:cont-deadroutes}

Each premise is calibrated against a live counterexample, and the counterexamples also
close the tempting short routes.

\emph{Superposition does not imply continuity.} Modal quantum theory
\cite{schumacherwestmoreland2012} has superposition over finite fields and no continuity
(and no probabilities). Spekkens's toy theory and the epistricted theories
\cite{spekkens2007, spekkens2016} are \emph{ignorance-based and discrete}: the
transformation group of a toy bit is finite, and the epistemic restriction is
integer-quantised (knowing one quadrature costs exactly one bit of the other). These
theories prove that ignorance in general does not force a continuum; the separating
premise is (P2$'$), which their finite groups violate maximally (a finite group has no
$n$-th roots beyond its exponent). Whether \emph{self-referential} ignorance can be
similarly quantised is Open Problem~\ref{op:quantised}.

\emph{Uncountability is not continuity, and the diagonal builds no completion.} Cantor's
argument is non-exhaustion, not construction: given a listing it exhibits an unlisted
element; it neither enumerates the gaps nor fills them, and uncountable totally
disconnected sets abound (Remark~\ref{rem:witnesses}). The diagonal's true role in this
programme is stated in \S\ref{app:cont-specker}.

\emph{Completeness is not connectedness.} $\mathbb{Q}_p$ (Remark~\ref{rem:witnesses})
--- hence the Archimedean selection of Lemma~\ref{lem:accumulation} is a premise doing
real work, and $p$-adic probability \cite{khrennikov-padic} is the consistent mathematics
of the road not taken.

\emph{Connected is not path-connected.} The solenoid --- excluded here only because
Cartan's theorem applies to groups of real matrices; that is, excluded by Layer~1.

\subsection{Super-tasks, Specker sequences, and what the diagonal does}
\label{app:cont-specker}

Specker \cite{specker1949} constructed a computable, monotone, bounded sequence of
rationals whose limit is not a computable real --- canonically
$\sigma = \sum_{n \in K} 2^{-n}$ with $K$ the halting set. This is the Thomson lamp
arithmetised: a dyadic super-task whose terminal value encodes halting information. It is
the exact formal witness for the situation (P1) legislates: a describable Cauchy sequence
whose limit the internal agent provably cannot realise. Three consequences. First, the
completion demanded by (P1) necessarily contains non-computable points --- the computable
reals are countable and Lebesgue-null --- so the continuity axiom is, on this reading,
the demand that the transformation group contain its \emph{non-computable} limits:
continuity is the topological face of the founding non-computability. Second, generic
super-tasks \emph{are} diagonalisations: Specker's $\sigma$ is a single act that is at
once a convergent super-task and an escape from every internal enumeration; the corpus
slogan identifying the two is cashed at the level of objects, not merely of licences.
Third, the record side has no canonical completion: infinite-time Turing machines
\cite{hamkinslewis2000} complete every $\omega$-super-task only by \emph{convention}
(the limsup rule assigns the lamp the value $1$; liminf would assign $0$) --- Benacerraf's
arbitrariness \cite{benacerraf1962} formalised --- whereas the amplitude-side fixed point
assigned in \S\ref{sec:fanout-boundary} is basis-free. Two ways to answer Thomson:
legislation or superposition; only the second is coordinate-free.

\subsection{Open problems}
\label{app:cont-openproblems}

\begin{openproblem}[Formalising descriptive closure]\label{op:fdp}
State (P1) as a principle about descriptions --- e.g.\ ``the theory's objects are the
finite-precision-specifiable classes, and its spaces the completions thereof'' --- and
either (a) derive it from the self-specification obstruction with no further physical
input, showing the closure $\overline{G_0}$ is the unique representative of $G_0$'s
operational-indistinguishability class definable without choice principles beyond
dependent choice; or (b) exhibit a countermodel: an operational theory whose
transformation set is describable, dense in its closure, and defensibly \emph{not}
closed. A solution to (a) turns Theorem~\ref{thm:conditional} into a derivation of
Axiom~5 from self-reference modulo Problems~\ref{op:stage-lattice}
and~\ref{op:transitivity}.
\end{openproblem}

\begin{openproblem}[The stage-lattice reduction: no records implies divisibility]
\label{op:stage-lattice}
Prove or refute: if some available transformation lacks an available $n$-th stage, then
its achievable-stage structure is a distinguished discrete invariant of the theory ---
a preferred-basis-like structure, hence a classical record --- contradicting the
pre-FANOUT regime's defining property. A proof reduces (P2$'$) to the regime's third
face, exactly as Problem~\ref{op:transitivity} reduces (P3), leaving
Problem~\ref{op:fdp} as the single bridge from self-referential ignorance proper. The
argument should have the shape of the preferred-basis proposition of
\S\ref{sec:fanout-boundary}: discreteness is what a record is.
\end{openproblem}

\begin{openproblem}[Is self-referential ignorance quantised or divisible?]
\label{op:quantised}
The epistricted theories show ignorance \emph{can} come in integer units; Chaitin's
$\Omega$ is maximally unknowable bit by bit. Show that \emph{pre-record} self-referential
ignorance admits no minimal unit --- e.g.\ because any unit would be a standard, a stable
copyable referent, hence a record --- so that bit-quantised ignorance ($\Omega$, FANOUT
outcomes) is post-record and divisible ignorance (phase) is pre-record, consistently
with the two-sided reading of the boundary. Alternatively, exhibit a discrete toy theory
whose epistemic restriction is genuinely self-referential; this would refute the
programme's strongest claim and is the honest falsifier.
\end{openproblem}

\begin{openproblem}[Transitivity from no records]\label{op:transitivity}
Make rigorous: orbits of the compact group $G$ on pure states are closed; a proper
invariant closed distinguished set of pure states defines a classical superselection
observable; a superselection observable is a record predating any FANOUT; hence in the
pre-FANOUT regime $G$ is transitive on pure states.
\end{openproblem}

\begin{openproblem}[The Lawvere unification]\label{op:lawvere}
Lawvere proved the diagonal fixed-point theorem in cartesian closed categories
\cite{lawvere1969} and, separately, characterised metric Cauchy completeness as
representability of limit-bimodules in enriched category theory \cite{lawvere1973}.
Conjecture: in a closed category with an amplitude object $A$ admitting a point-surjective
self-describer $e : X \to A^{X}$ (the categorical form of inside/outside equivalence),
the enriched homs into $A$ must be Cauchy-complete --- an unrepresented describable
limit-bimodule yields a diagonalisation witness against point-surjectivity (``the
description refers to a transformation that does not exist''). A proof would derive (P1)
inside one mathematician's two theorems and would constitute the sought result in full.
\end{openproblem}

\subsection{Relation to existing work}
\label{app:cont-related}

Every individual mathematical ingredient above is classical: Cartan's closed-subgroup
theorem; the Baire argument for countable compact groups; Ostrowski's classification;
Specker's sequence; the automatic-continuity literature; the structure theory of compact
groups \cite{hofmannmorris}. What we have not found anywhere is the \emph{question}: no
reconstruction of quantum theory derives its continuity axiom from anything --- Hardy
postulates Axiom~5 \cite{hardy2001}; Schack's four-axiom variant retains it
\cite{schack2003}; continuous reversibility is likewise postulated by Dakić--Brukner and
Masanes--M\"uller \cite{dakicbrukner2011, masanesmueller2011}, and the
operational-probabilistic framework of Chiribella--D'Ariano--Perinotti \emph{assumes}
topological closure of its procedure sets \cite{cdp2011} --- the assumption (P1) makes
explicit and seeks to ground. The nearest programmes run adjacent arrows: computable
analysis and domain theory prove computable-implies-continuous, but in a semantic
category rather than an operational state space \cite{weihrauch2000, abramskyjung};
Gisin's finite-information physics argues from finite information \emph{against} the
completed continuum \cite{gisin2019, gisin2021} --- the reverse of the present arrow,
which holds that finite description forces the completion \emph{as the description
space}, importing the non-computable as superposition rather than expelling it; Svozil
connects the lamp, diagonalisation, and superposition as fixed points, at the level of
analogy \cite{svozil2009}; and the finite-precision Kochen--Specker debate
\cite{meyer1999, kent1999, cliftonkent2000} is exactly a dispute about whether the
completion carries physics its dense describable skeleton lacks --- this framework
answers yes, and locates the content in self-reference. We therefore offer
Problems~\ref{op:fdp}--\ref{op:lawvere} in the expectation that either they are new, or
a reader will point us to their solutions; both outcomes advance the programme, and the
problems are stated so that a specialist can engage each without adopting the framework
whole.

% ---------------------------------------------------------------------
% BIBLIOGRAPHY STUBS (informal; VERIFY all at final drafting)
% hardy2001: L. Hardy, Quantum Theory From Five Reasonable Axioms,
%   arXiv:quant-ph/0101012 (2001).
% schack2003: R. Schack, Quantum theory from four of Hardy's axioms,
%   Found. Phys. 33, 1461 (2003); arXiv:quant-ph/0210017.
% dakicbrukner2011: B. Dakic, C. Brukner, Quantum theory and beyond...,
%   in Deep Beauty (ed. Halvorson), CUP (2011); arXiv:0911.0695.
% masanesmueller2011: L. Masanes, M. Mueller, A derivation of quantum
%   theory from physical requirements, New J. Phys. 13, 063001 (2011).
% cdp2011: G. Chiribella, G. M. D'Ariano, P. Perinotti, Informational
%   derivation of quantum theory, PRA 84, 012311 (2011).
% lawvere1969: F. W. Lawvere, Diagonal arguments and cartesian closed
%   categories, LNM 92 (1969); reprinted TAC.
% lawvere1973: F. W. Lawvere, Metric spaces, generalized logic, and
%   closed categories, Rend. Sem. Mat. Fis. Milano 43 (1973).
% specker1949: E. Specker, Nicht konstruktiv beweisbare Saetze der
%   Analysis, J. Symbolic Logic 14, 145-158 (1949).
% weihrauch2000: K. Weihrauch, Computable Analysis, Springer (2000).
% abramskyjung: S. Abramsky, A. Jung, Domain theory, in Handbook of
%   Logic in Computer Science III (1994).
% smyth: M. B. Smyth, Topology, in Handbook of Logic in Computer
%   Science I (1992) -- opens as observable properties.
% rosendal2009: C. Rosendal, Automatic continuity of group
%   homomorphisms, Bull. Symbolic Logic 15 (2009).
% vanderwaerden1933: B. L. van der Waerden, Stetigkeitssaetze fuer
%   halbeinfache Liesche Gruppen, Math. Z. 36 (1933).
% tsankov2013: T. Tsankov, Automatic continuity for the unitary group,
%   Proc. AMS / arXiv:1109.1200 (VERIFY venue).
% gleason1952: A. M. Gleason, Groups without small subgroups,
%   Ann. Math. 56 (1952).
% gleason1957: A. M. Gleason, Measures on the closed subspaces of a
%   Hilbert space, J. Math. Mech. 6 (1957).
% montgomeryzippin: D. Montgomery, L. Zippin, Topological
%   Transformation Groups, Interscience (1955).
% hofmannmorris: K. H. Hofmann, S. A. Morris, The Structure of Compact
%   Groups, de Gruyter -- compact abelian: connected iff divisible
%   (VERIFY theorem number).
% ostrowski1916: A. Ostrowski, Ueber einige Loesungen der
%   Funktionalgleichung..., Acta Math. 41 (1916).
% kolmogorov1933: A. N. Kolmogorov, Grundbegriffe der
%   Wahrscheinlichkeitsrechnung (1933), axiom V.
% definetti: B. de Finetti, Theory of Probability, Wiley (1974) --
%   finite additivity.
% popper1950: K. Popper, Indeterminism in quantum physics and in
%   classical physics I & II, BJPS 1 (1950). [project PDFs]
% thomson1954: J. F. Thomson, Tasks and super-tasks, Analysis 15 (1954).
% benacerraf1962: P. Benacerraf, Tasks, super-tasks, and the modern
%   Eleatics, J. Phil. 59 (1962).
% hamkinslewis2000: J. D. Hamkins, A. Lewis, Infinite time Turing
%   machines, J. Symbolic Logic 65 (2000).
% spekkens2007: R. W. Spekkens, Evidence for the epistemic view of
%   quantum states: a toy theory, PRA 75, 032110 (2007).
% spekkens2016: R. W. Spekkens, Quasi-quantization: classical
%   statistical theories with an epistemic restriction (2016).
% schumacherwestmoreland2012: B. Schumacher, M. Westmoreland, Modal
%   quantum theory, Found. Phys. 42 (2012) (VERIFY year).
% khrennikov-padic: A. Khrennikov, p-adic probability and statistics /
%   Non-Archimedean analysis (VERIFY best citation).
% meyer1999: D. A. Meyer, Finite precision measurement nullifies the
%   Kochen-Specker theorem, PRL 83, 3751 (1999).
% kent1999: A. Kent, Noncontextual hidden variables and physical
%   measurements, PRL 83, 3755 (1999).
% cliftonkent2000: R. Clifton, A. Kent, Simulating quantum mechanics by
%   non-contextual hidden variables, Proc. R. Soc. A 456 (2000).
% gisin2019: N. Gisin, Indeterminism in physics, classical chaos and
%   Bohmian mechanics..., arXiv:1803.06824 / Erkenntnis (VERIFY).
% gisin2021: N. Gisin, Indeterminism in physics and intuitionistic
%   mathematics, Synthese 199 (2021).
% svozil2009: K. Svozil, arXiv:0904.1649 (RECORDED in corpus).
% =====================================================================
